\documentclass[12pt]{article}

\usepackage[a4paper,margin=1in]{geometry}
\usepackage{amsmath}
\usepackage{setspace}
\usepackage{enumitem}
\usepackage[version=4]{mhchem}

\title{\textbf{Preparation of 0.1 M Acetate Buffer (pH 5.0)}}
\author{}
\date{}

\begin{document}
	
	\maketitle
	
	\section*{Aim}
	
	To prepare 100 mL of 0.1 M acetate buffer of pH 5.0 using 0.1 M sodium acetate and 0.1 M acetic acid solutions.
	
	\section*{Required Solutions}
	
	\begin{itemize}
		\item 0.1 M Acetic acid solution (\(\mathrm{CH_3COOH}\))
		\item 0.1 M Sodium acetate solution (\(\mathrm{CH_3COONa}\))
		\item Distilled water (\(\mathrm{H_2O}\))
	\end{itemize}
	
	\section*{Apparatus}
	
	\begin{itemize}
		\item 100 mL volumetric flask
		\item measuring cylinders
		\item Beaker
		\item Glass rod
		\item pH meter
	\end{itemize}
	
	\section*{Theory}
	
	An acetate buffer consists of a weak acid (acetic acid, \(\mathrm{CH_3COOH}\)) and its conjugate base (sodium acetate, \(\mathrm{CH_3COONa}\)). It resists changes in pH when small amounts of acid or base are added. The pH of the buffer is governed by the Henderson--Hasselbalch equation:
	
	\[
	\mathrm{pH}=\mathrm{p}K_a+\log\left(\frac{[\mathrm{CH_3COO^-}]}{[\mathrm{CH_3COOH}]}\right)
	\]
	
	For acetic acid,
	
	\[
	\mathrm{p}K_a=4.76
	\]
	
	For a buffer of pH 5.0,
	
	\[
	5.0=4.76+\log\left(\frac{[\mathrm{CH_3COO^-}]}{[\mathrm{CH_3COOH}]}\right)
	\]
	
	Therefore,
	
	\[
	\frac{[\mathrm{CH_3COO^-}]}{[\mathrm{CH_3COOH}]}
	=10^{(5.0-4.76)}
	=10^{0.24}
	\approx1.74
	\]
	
	Since both stock solutions have the same concentration (0.1 M), the concentration ratio is equal to the volume ratio.
	
	Thus,
	
	\[
	\frac{V_{\mathrm{CH_3COONa}}}
	{V_{\mathrm{CH_3COOH}}}
	=1.74
	\]
	
	For a total volume of 100 mL,
	
	\[
	V_{\mathrm{CH_3COONa}}
	=
	\frac{1.74}{1+1.74}\times100
	=
	63.5~\mathrm{mL}
	\]
	
	\[
	V_{\mathrm{CH_3COOH}}
	=
	\frac{1}{1+1.74}\times100
	=
	36.5~\mathrm{mL}
	\]
	
	Hence, mixing 63.5 mL of 0.1 M sodium acetate with 36.5 mL of 0.1 M acetic acid produces a 0.1 M acetate buffer of pH 5.0.
	
	\section*{Procedure}
	
	\begin{enumerate}
		\item Clean a 100 mL volumetric flask and rinse it with distilled water.
		\item Measure 63.5 mL of 0.1 M sodium acetate (\(\mathrm{CH_3COONa}\)) using measuring cylinder and add into the volumetric flask.
		\item Add 36.5 mL of 0.1 M acetic acid (\(\mathrm{CH_3COOH}\)) into above solution.
		\item Mix the solution thoroughly with constant stirring.
		\item Measure the pH using a calibrated pH meter.
		\item If necessary, adjust the pH:
		\begin{itemize}
			\item Add a few drops of 0.1 M acetic acid to decrease the pH.
			\item Add a few drops of 0.1 M sodium acetate to increase the pH.
		\end{itemize}
		\item If total volume is slightly less than 100 mL after pH adjustment, make up the volume to mark with distilled water.
	\end{enumerate}
	
	\section*{Observation}
	
	The prepared solution was clear and colourless. The measured pH was approximately 5.0.
	
	\section*{Result}
	
	A 100 mL sample of 0.1 M acetate buffer of pH 5.0 was successfully prepared by mixing 63.5 mL of 0.1 M sodium acetate solution with 36.5 mL of 0.1 M acetic acid solution.
	
	\section*{Precautions}
	
	\begin{enumerate}
		\item Wash the Apparatus carefully.
		\item pH meter must be calibrated before use.
		\item Handle Acid solution with care.
\end{enumerate}

	
\end{document}
